\section*{Chapter \ref{ch:rs:event-driven-backtests} --- Event-Driven Backtests}

\begin{solution}{exo:rs:event-driven-backtests:1}
\omterm{def:rs:event-driven-backtests:fill}{Touch fill}: the low reached 99.00, so yes, in full at 99.00. Penetration by one tick: the low would have to reach 98.99; it did not, so no.
\end{solution}

\begin{solution}{exo:rs:event-driven-backtests:2}
2.5\% of each \omterm{def:rs:market-data-for-research:bar}{bar}: 1\,000, 250 and 1\,500 shares, 2\,750 in all, leaving 2\,250 working after the third \omterm{def:rs:market-data-for-research:bar}{bar}.
\end{solution}

\begin{solution}{exo:rs:event-driven-backtests:3}
pending (sent) $\to$ working (acknowledged) $\to$ partial (100 filled, 200 remaining) $\to$ cancelled (the 200 withdrawn when the cancellation takes effect);
the 100 filled shares stay filled.
\end{solution}

\begin{solution}{exo:rs:event-driven-backtests:4}
A trade through the order's price means the whole queue at that price was consumed, which happens when the flow on the other side is heavy and
persistent, often informed (chapter~9); the price then keeps moving. A touch can happen with a queue barely dented and the price turning back. The
\omterm{def:rs:event-driven-backtests:fill}{penetration fill} therefore selects the fills a real order at the back of the queue would get, and those are the adverse ones: on the synthetic days,
$-2.0$ ticks of mark-out against $-0.16$.
\end{solution}

\begin{solution}{exo:rs:event-driven-backtests:5}
Every order is sent at a \omterm{def:rs:market-data-for-research:bar}{bar}'s close, becomes live 5 seconds into the next \omterm{def:rs:market-data-for-research:bar}{bar} and is cancelled 5 seconds into the one after; it is never live for a whole
\omterm{def:rs:market-data-for-research:bar}{bar}. A strategy that leaves its orders alone while the price is unchanged keeps them live for many whole \omterm{def:rs:market-data-for-research:bar}{bars}, and the conservative policy can fill them.
\end{solution}

\begin{solution}{exo:rs:event-driven-backtests:6}
At the split's effective time: position to 200 shares, the order to 200 shares at \$60. Without it, the order still asks \$120 for 100 shares in a market
now around \$60: it never fills (the book is unhedged against its intended exit), or, for a buy limit left at the old price, it fills at once at twice the
market.
\end{solution}

\begin{solution}{exo:rs:event-driven-backtests:7}
\texttt{rs\_evbt.at\_close}: under the \omterm{def:rs:event-driven-backtests:fill}{touch fill} 95\% of orders fill, with a mark-out of $-0.17$ ticks and a loss of \$188 a day; under the \omterm{def:rs:event-driven-backtests:fill}{penetration fill}
70\% fill, with $-1.16$ ticks and a loss of \$506 a day. Orders at the price that just traded are touched almost always, which says nothing about whether a
real order at the back of its queue would fill.
\end{solution}

\begin{solution}{exo:rs:event-driven-backtests:8}
The \omterm{def:rs:event-driven-backtests:fill}{fill model} assumes the order was first in the queue at its price, and fills it on every touch, including all the touches where the price turned
back, the fills a real order would not get; it ignores the adverse fills' selection. Rerun with \omterm{def:rs:event-driven-backtests:fill}{penetration fills} and a volume cap, check the mark-outs,
and confirm at level~3 or with live fills before believing it.
\end{solution}

\begin{solution}{pb:rs:event-driven-backtests:1}
\begin{enumerate}
\item Fills against the \omterm{def:rs:market-data-for-research:bar}{bar} first (the orders were working during it), then the mark at its close, then the strategy's reaction; so the strategy sees
the fills and the portfolio it actually had, and its new orders act only on later \omterm{def:rs:market-data-for-research:bar}{bars}.
\item When its acknowledgement event arrives, after the order-entry latency; and under the conservative policy only in \omterm{def:rs:market-data-for-research:bar}{bars} that begin after that.
\item A total order on events (time, kind, sequence number), seeded randomness passed in explicitly, no wall-clock dependence.
\item The \texttt{BacktestResult} of chapter~16, the orders with their lifecycles, and the fills.
\item Touch: 523 lots, 70\%, $+\$17$, $-0.16$ ticks. Penetration: 290 lots, 41\%, $-\$584$, $-2.00$ ticks. Capped at 2.5\%: 388 lots, 55\%, $+\$20$,
$-0.18$ ticks.
\item Touch: first in the queue. Penetration: last in the queue. Capped: a fixed share of the volume, wherever the order stands.
\item Its fills are only those where the price traded through, which are followed by moves against the position; the \omterm{def:rs:event-driven-backtests:fill}{touch fill} adds the benign fills
where the price turned back.
\item Hold one lot long after a down \omterm{def:rs:market-data-for-research:bar}{bar} and one short after an up \omterm{def:rs:market-data-for-research:bar}{bar}, traded at the close with half a tick of cost: $-\$215$ a day.
\item No fills at all: no order is live for a whole \omterm{def:rs:market-data-for-research:bar}{bar}.
\item Touch: 585 lots, $-\$41$; penetration: 372 lots, $-\$776$; capped: 440 lots, $-\$190$, with its mark-out falling to $-0.87$ ticks.
\item That the answer depends on when within the \omterm{def:rs:market-data-for-research:bar}{bar} the price reached the orders, which \omterm{def:rs:market-data-for-research:bar}{bars} do not record: the strategy must be judged at level~3.
\item Keep orders working while the price is unchanged, cancel only when it moves, and quote further from the price.
\item \textbf{Named result.} With no latency: touch 70\% of orders filled and $+\$17$ a day, penetration 41\% and $-\$584$, capped 55\% and $+\$20$; with a
5-second latency: nothing under the conservative policy, and $-\$41$, $-\$776$ and $-\$190$ under the optimistic one. The \omterm{def:rs:event-driven-backtests:fill}{fill model} moves the answer by \$800
a day; the strategy's own edge is within its noise.
\item The penetration result, as the pessimistic bound, with the statement that the touch result assumes first place in every queue and that the level-2
results bracket the answer without locating it.
\item Our own orders in the queue at their prices, every trade at those prices, and latency on both the order and the data side: fills that happen or not
for the reasons they would live.
\item At most 2.5\% of each \omterm{def:rs:market-data-for-research:bar}{bar} fills; a strategy that needs 5\% leaves half its orders working and trades later, at other prices.
\item Splits, dividends (for the cash and the price gap), half-days, auctions and halts.
\item Because a real order joins the back of a queue that may hold thousands of shares at its price; to fill 70\% of the time, the price would have to trade
through that queue in most minutes.
\item Each order's send, acknowledgement, queue estimate, fills and cancellations with timestamps, and the market's trades at its price; the ratio of live
fills to \omterm{def:rs:event-driven-backtests:fill}{touch fills} calibrates the model.
\item The \omterm{def:rs:vectorised-backtests:backtest}{backtest}'s assumption about where its orders stood in the queue, made explicit and replaceable.
\end{enumerate}
\end{solution}

\begin{solution}{iq:rs:event-driven-backtests:1}
A priority queue of time-stamped events (market data, order acknowledgements, fills, cancellations, calendar and corporate events), a dispatcher that pops
them in order and advances the \omterm{def:rs:event-driven-backtests:evbt}{simulated clock}, handlers for the strategy, the order manager (lifecycle), the \omterm{def:rs:event-driven-backtests:fill}{fill model} and the portfolio, all of which may
only schedule events in the future; and a recorder that produces the result and the audit trail.
\end{solution}

\begin{solution}{iq:rs:event-driven-backtests:2}
\omterm{def:rs:event-driven-backtests:fill}{Touch fills} (optimistic: first in the queue), \omterm{def:rs:event-driven-backtests:fill}{penetration fills} (pessimistic: last in the queue), volume-capped fills (a share of the \omterm{def:rs:market-data-for-research:bar}{bar}'s volume),
probabilistic fills calibrated on live data; each needs a policy for orders live for part of a \omterm{def:rs:market-data-for-research:bar}{bar}. Biases: touch overstates fills and understates adverse
selection; penetration understates fills; caps ignore where the volume traded.
\end{solution}

\begin{solution}{iq:rs:event-driven-backtests:3}
Order events totally (time, then kind, then a sequence number), pass seeded generators explicitly, avoid iteration over unordered containers and any use of the
wall clock, and test that two runs produce identical results.
\end{solution}

\begin{solution}{iq:rs:event-driven-backtests:4}
The tendency of passive orders to be filled when the price is about to move against them, because the flow that fills them is informed or persistent.
Measure the mark-out of each fill: the mid a few seconds or \omterm{def:rs:market-data-for-research:bar}{bars} later minus the fill price, in the fill's direction, averaged by \omterm{def:rs:event-driven-backtests:fill}{fill model}, time and size.
\end{solution}

\begin{solution}{iq:rs:event-driven-backtests:5}
As an event at its effective time: multiply positions and working quantities by the ratio, divide limit and stop prices by it, keep the \omterm{def:rs:market-data-for-research:bar}{bars} as traded, and log
the adjustment; the same for reverse splits, and a check that capital is unchanged by the event.
\end{solution}

\begin{solution}{iq:rs:event-driven-backtests:6}
When the result depends on assumptions the \omterm{def:rs:market-data-for-research:bar}{bars} cannot settle: run it under touch and \omterm{def:rs:event-driven-backtests:fill}{penetration fills}, and under conservative and optimistic live-in-bar
policies; if the answers bracket zero or differ by more than the edge, it needs level~3 (queue and latency) or live evidence.
\end{solution}
